\chapter{Introdução}
\label{chapter:introduction}

No mundo digital, copiar informação é trivial; alterá-la também. Para que uma
moeda digital seja aceite, os seus utilizadores precisam de confiar em duas
coisas: a quantidade total de moeda obedece às regras anunciadas e cada unidade
só pode ser gasta uma vez. Quem recebe um pagamento tem, portanto, de poder
verificar que as moedas não são contrafeitas nem foram já entregues a outra
pessoa. Se não quisermos confiar essa tarefa a uma autoridade central, a oferta
monetária e a história completa das transacções têm de ser verificáveis pelo
público.

As cripto-moedas registam transacções numa lista de blocos
(\emph{blockchain}), uma base de dados distribuída cuja história é muito
difícil de alterar sem que a rede o detecte \cite{Nakamoto_bitcoin}. Muitas
listas de blocos expõem os intervenientes e os montantes para tornar simples a
verificação colectiva. O preço é evidente: perde-se privacidade e, com ela,
fungibilidade.\footnote{Um bem é \textbf{fungível} quando uma unidade pode
substituir outra no uso ou no cumprimento de um contracto
\cite{mises-org-fungible}. Numa lista de blocos transparente como a de Bitcoin,
as moedas podem ser distinguidas pela sua história. Se essa história incluir
actores considerados suspeitos, certas moedas podem ser tratadas como
``manchadas'' \cite{bitcoin-big-bang-taint}; inversamente, houve quem pagasse
um prémio por moedas recém-mineradas, sem história anterior
\cite{new-bitcoin-premium}.}

Um utilizador de Bitcoin pode tentar ofuscar o percurso dos fundos com
endereços intermédios \cite{DBLP:journals/corr/NarayananM17}. Ainda assim,
ferramentas de análise conseguem muitas vezes estabelecer ligações entre
remetentes e destinatários
\cite{DBLP:journals/corr/ShenTuY15b, DK-police-tracing-btc,
Andrew-Cox-Sandia, chainalysis-2020-report}.

Monero adopta outra estratégia. Endereços de utilização única escondem o
destinatário na lista de blocos; assinaturas em anel tornam ambígua a saída que
foi efectivamente gasta; e compromissos criptográficos ocultam os montantes sem
impedir a verificação da contabilidade. O resultado é uma cripto-moeda com um
elevado grau de privacidade e fungibilidade. Estas propriedades não eliminam a
informação produzida pelo comportamento do utilizador ou disponível fora da
lista de blocos, que pode permitir alguma análise.\footnote{Para um exemplo de heurísticas
aplicadas às assinaturas em anel, veja-se \cite{monero-ring-heuristics-ryo}.}

\section{Objectivos}
\label{sec:goals}

Monero foi lançado em 2014 e tem uma história continuada de desenvolvimento
\cite{bitmonero-launched, monero-history}. Quando a segunda edição deste livro
foi preparada, a comunidade e a adopção da moeda continuavam a aumentar
\cite{justin-defcon-2019-community-growth}.\footnote{Naquele período, Monero
ocupava o 14.º lugar por capitalização de mercado em 14 de
Junho de 2018 e o 12.º em 5 de Janeiro de 2020; veja-se
\url{https://coinmarketcap.com/}. Estes valores não são uma medida técnica nem
uma afirmação sobre a posição actual.}

Documentar o sistema é difícil. Existe documentação prática no projecto
\url{https://monerodocs.org/} e em \emph{Mastering Monero}
\cite{mastering-monero}, mas partes importantes da construção teórica foram
publicadas em textos incompletos, não revistos por pares ou entretanto
ultrapassados. Para saber o que uma versão do programa faz, o código-fonte é a
referência decisiva; para saber por que uma construção criptográfica é segura,
continuam a ser indispensáveis os artigos e as provas que a fundamentam.
\footnote{A série \emph{Monero Building Blocks}, de Seguias
\cite{monero-building-blocks}, trata com cuidado as provas de segurança dos
esquemas de assinatura. Tal como a primeira edição de \emph{Zero a Monero}
\cite{ztm-1}, concentra-se no protocolo v7.}

Para leitores sem formação matemática, os fundamentos da criptografia de
curva elíptica podem exigir conceitos preliminares que outras descrições
pressupõem.
Uma explicação anterior do funcionamento de Monero
\cite{MRL-0003-about-monero}, por exemplo, não desenvolvia esses fundamentos,
era incompleta e ficou obsoleta. Este livro introduz os conceitos necessários,
revê os algoritmos criptográficos e reúne uma descrição construtiva, passo a
passo, do funcionamento interno de Monero.

\section{Âmbito técnico desta edição}
\label{sec:technical-scope}

A segunda edição original descreve o protocolo v12 e o programa Monero
v0.15.x.x. A revisão portuguesa toma esse texto como ponto de partida, mas
confere o estado corrente contra a rede principal e a revisão do código
identificada no resumo inicial.\footnote{Foram conferidos, entre
outros, \path{README.md}, \path{src/hardforks/hardforks.cpp},
\path{src/cryptonote_core/blockchain.cpp},
\path{src/cryptonote_protocol/levin_notify.cpp} e o directório
\path{src/fcmp_pp/}. O repositório completo está em
\url{https://github.com/monero-project/monero}.}

Esta distinção evita três confusões frequentes:

\begin{itemize}
    \item uma \emph{regra de consenso activa} determina que blocos e
    transacções a rede principal aceita;
    \item um \emph{comportamento implementado} pode pertencer à carteira ou à
    rede entre nós sem fazer parte do consenso;
    \item \emph{trabalho experimental} pode existir no código antes de estar
    completo, estável ou sequer programado para activação.
\end{itemize}

Na v16, as novas transacções usam CLSAG e Bulletproof+, anéis fixos de 16
membros e etiquetas de vista. RandomX é a prova de trabalho activa desde a v12.
Dandelion++ é comportamento implementado para a difusão de transacções na rede
pública, não uma regra de consenso. O ramo consultado contém componentes de
FCMP++ e tipos associados a Carrot, mas não programa uma versão posterior à
v16 na rede principal; são, por isso, trabalho experimental nesta edição.

O ``protocolo'' é o conjunto de regras contra o qual cada novo bloco é testado
antes de entrar na lista. Inclui as regras das transacções, mas não se confunde
com o campo de versão da própria transacção. Este último continua a identificar
RingCT v2, embora várias regras concretas e o formato RingCT tenham mudado de
facto ao longo das versões de consenso.

A actualização desta edição abrange todos os capítulos. As descrições do
protocolo activo referem-se à v16; os formatos e esquemas de épocas anteriores
são identificados expressamente como históricos quando ajudam a interpretar
dados antigos ou a acompanhar uma derivação. A primeira edição, por sua vez,
corresponde ao protocolo v7 e ao programa v0.12.x.x \cite{ztm-1}.

O código beneficia de revisão pública, mas isso não o torna infalível. O leitor
é convidado a conferir as explicações nas fontes. Uma vulnerabilidade séria
deve ser comunicada de forma responsável em
\url{https://hackerone.com/monero}; uma correcção não sensível pode ser proposta
no repositório do projecto. As propostas Triptych \cite{triptych-preprint},
RingCT 3.0 \cite{ringct3-preprint}, Omniring \cite{omniring-paper} e Lelantus
\cite{lelantus-preprint} pertencem ao contexto de investigação da segunda
edição, não ao conjunto de regras activas aqui descrito.

\section{Aos leitores}

Prevemos que muitos leitores cheguem a este relatório com pouco ou nenhum
contacto com matemática discreta, estruturas algébricas, criptografia ou listas
de blocos. Procurámos dar pormenor suficiente para que um leitor atento possa
aprender sem depender constantemente de pesquisa externa.\footnote{Uma
introdução extensa à criptografia aplicada encontra-se em
\cite{applied-cryptography-textbook}.}

Alguns pormenores matemáticos foram omitidos ou remetidos para notas de rodapé
quando prejudicariam a clareza. Também não seguimos cada detalhe de
implementação quando ele não é necessário para compreender o mecanismo. O
objectivo é combinar criptografia matemática e implementação: apresentar as
hipóteses formalmente e mostrar de forma concreta como os componentes se
relacionam.
\footnote{Certas notas de rodapé antecipam capítulos posteriores. Foram
pensadas para ganhar utilidade numa segunda leitura, quando os detalhes de
implementação já têm onde assentar.}

\section{Origens da cripto-moeda Monero}

Monero, inicialmente chamado BitMonero, foi criado em Abril de 2014 a partir da
prova de conceito CryptoNote \cite{bitmonero-launched}. \emph{Monero} significa
``moeda'' em esperanto; o plural é \emph{moneroj}.

CryptoNote foi desenvolvido por várias pessoas. O texto de referência que o
descreve foi publicado sob o pseudónimo Nicolas van Saberhagen, em Outubro de
2013 \cite{cryptoNoteWhitePaper}. Propunha endereços de utilização única para
ocultar o destinatário e assinaturas em anel para tornar ambíguo o remetente.
Desde então, Monero reforçou a privacidade com montantes confidenciais,
inspirados, entre outros, no trabalho de Greg Maxwell
\cite{Signatures2015BorromeanRS} e integrados em RingCT segundo as recomendações
de Shen Noether \cite{MRL-0005-ringct}. As provas Bulletproof tornaram depois
essa ocultação muito mais eficiente \cite{Bulletproofs_paper}; CLSAG reduziu o
custo das assinaturas em anel \cite{MRL-0011-CLSAG}, e Bulletproof+ sucedeu às
provas Bulletproof nas transacções activas.

\subsection{Parte 1: ``Essenciais''}

Começamos pelas ferramentas criptográficas necessárias. O capítulo
\ref{chapter:basicConcepts} desenvolve aritmética modular, curvas elípticas,
chaves, assinaturas de Schnorr e a notação usada no resto do livro. O capítulo
\ref{chapter:advanced-schnorr} prolonga essa construção até às assinaturas em
anel. No capítulo \ref{chapter:addresses}, vemos como os endereços controlam a
recepção de fundos e por que uma saída não revela publicamente o destinatário.

O capítulo \ref{chapter:pedersen-commitments} introduz compromissos e provas de
intervalo para ocultar montantes. Em conjunto, o capítulo
\ref{chapter:transactions} constrói uma transacção RingCT. O capítulo
\ref{chapter:blockchain} abre por fim a lista de blocos: mineração, dificuldade,
emissão, peso e taxas. O capítulo \ref{chapter:randomx} segue uma tentativa de
RandomX desde a semente de época até ao hash de prova de trabalho e explica
por que programas aleatórios e memória abundante procuram reduzir a vantagem
de máquinas especializadas. O capítulo \ref{chapter:dandelion} acompanha
depois a vida de um nó: Dandelion++ afasta a origem de uma transacção do ponto
onde começa a difusão pública, enquanto a poda e a sincronização podada
distribuem pela rede a conservação da história sem as confundir com
anonimato. O capítulo \ref{chapter:fcmp-plus-plus} analisa a
alteração FCMP++ ainda em
desenvolvimento: parte das provas de Schnorr e das árvores de curvas para
explicar como uma entrada poderá pertencer privadamente ao conjunto de saídas da cadeia
inteira. Como fecho dos ``Essenciais'', o capítulo
\ref{chapter:dlog-relations} pergunta o que aconteceria se alguém descobrisse
as relações supostamente desconhecidas entre os geradores: quando se perde só
privacidade, quando se torna possível o gasto duplo e quando aparece inflação.
Esta ordem fornece a base criptográfica para a parte seguinte; não transforma
a proposta FCMP++ em regra activa da versão 16.

\subsection{Parte 2: ``Extensões''}

Uma cripto-moeda é mais do que as suas regras de consenso. Esta parte explora
aplicações e propostas, várias das quais não estão implementadas ou dependem de
funcionalidades experimentais. Uma futura alteração do formato das transacções
pode torná-las impraticáveis; cada capítulo deve, portanto, ser lido com o seu
estatuto técnico em mente.

O capítulo \ref{chapter:tx-knowledge-proofs} mostra que informação sobre uma
transacção pode ser demonstrada a um observador. O capítulo
\ref{chapter:multisignatures} descreve a abordagem de multi-assinaturas da
carteira; na implementação consultada, esta funcionalidade continua marcada
como experimental e desactivada por omissão. O capítulo
\ref{chapter:escrowed-market} usa multi-assinaturas para desenhar mercados com
garantia. O capítulo \ref{chapter:txtangle} apresenta TxTangle, uma proposta
original e descentralizada para reunir as transacções de várias pessoas numa
só; o capítulo conserva a sua álgebra, mas deixa claro que o protocolo não foi
implementado no Monero.

\subsection{Conteúdo adicional}

O apêndice \ref{appendix:RCTTypeBulletproof2} disseca uma transacção
\texttt{RCTTypeBulletproof2}; é um exemplo histórico, não o formato que a v16
aceita para novas transacções. O apêndice \ref{appendix:block-content} explica
a estrutura de um bloco, incluindo o cabeçalho e a transacção de mineiro. O
apêndice \ref{appendix:genesis-block} fecha o percurso com o bloco génese de
Monero. O apêndice \ref{appendix:nota} fixa a notação vectorial e deriva o
termo \(\delta(y,z)\) das Bulletproofs sem saltos algébricos. Em conjunto,
ligam a teoria a objectos concretos da lista de blocos e às identidades que os
verificam.

As notas de margem \marginnote{isto é útil!} apontam para locais relevantes no
código-fonte. Os caminhos modernos foram conferidos na revisão indicada no
resumo; quando se conserva um caminho da v0.15.x.x para acompanhar uma
construção histórica, o contexto assinala essa época. A história completa
continua disponível em Git. Uma nota contém normalmente um caminho, como
\path{src/ringct/rctOps.cpp}, e por vezes uma função, como
\texttt{ecdhEncode()}. Um hífen num caminho estreito pode ser apenas uma quebra
tipográfica; os qualificadores de \emph{namespace}, como
\texttt{Blockchain::}, são geralmente omitidos.

\section{Aviso}

Os esquemas de assinatura, as aplicações de curvas elípticas e os detalhes da
implementação aqui apresentados são descritivos. Quem pretenda construir um
sistema real deve consultar as fontes primárias, as especificações técnicas e
o código correspondente à versão que vai usar.

Um esquema de assinatura precisa de uma prova de segurança bem definida e de
revisão competente. Vale aqui a velha regra: não invente a sua própria
criptografia. Código que implemente primitivas criptográficas deve ser revisto
por especialistas e acumular um historial longo e fiável. As contribuições
originais deste documento podem não ter recebido a mesma revisão nem testes;
leia-as com o cuidado devido.

\section{A história de Zero a Monero}

\emph{Zero a Monero} começou como expansão da tese de mestrado de Kurt Alonso,
``Monero --- Privacy in the Blockchain'' \cite{kurt-original}, publicada em
Maio de 2018. A primeira edição do livro apareceu em Junho desse ano
\cite{ztm-1}.

Para a segunda edição, os autores melhoraram a introdução às assinaturas em
anel (capítulo \ref{chapter:advanced-schnorr}) e reorganizaram a explicação das
transacções, acrescentando o capítulo de endereços
(\ref{chapter:addresses}). Modernizaram a comunicação dos montantes de saída
(secção \ref{sec:pedersen_monero}), substituíram as provas Borromean por
Bulletproof no percurso corrente da época (secção \ref{sec:range_proofs}),
retiraram \texttt{RCTTypeFull} do formato corrente e actualizaram o peso
dinâmico dos blocos e as taxas dos mineiros.

Foram ainda estudadas provas sobre transacções (capítulo
\ref{chapter:tx-knowledge-proofs}), multi-assinaturas (capítulo
\ref{chapter:multisignatures}), mercados com garantia (capítulo
\ref{chapter:escrowed-market}) e a proposta TxTangle (capítulo
\ref{chapter:txtangle}). Essa edição ficou alinhada com o protocolo v12 e o
programa v0.15.x.x.\footnote{O código-fonte \LaTeX{} das duas edições está em
\url{https://github.com/UkoeHB/Monero-RCT-report}; a primeira edição encontra-se
no ramo \texttt{ztm1}.}

\section{Reconhecimentos}
\label{sec:acknowledgements}

Texto original do autor ``koe''.

Este relatório não existiria sem a tese de mestrado de Kurt Alonso
\cite{kurt-original}, a quem o autor deve profunda gratidão. Brandon ``Surae
Noether'' Goodell e o investigador conhecido pelo pseudónimo ``Sarang Noether'',
do Monero Research Lab, foram fontes constantes de conhecimento durante as duas
edições. ``moneromooo'', um dos programadores mais prolíficos do projecto,
identificou inúmeras vezes as partes relevantes do código. Agradecemos também a
todos os contribuidores que responderam a perguntas e aos leitores que enviaram
correcções e palavras de encorajamento.
